\section{Scope and limitations}
\label{sec:limitations}

The results have a deliberately narrow interpretation.  \Profile{} proves the
atomic-v3 relation in Definition~\ref{def:spend-relation}: one input note
is replaced by one output note along the same depth-20 private path.  It does
not provide multiple inputs or outputs, a public append, recursion, proof
aggregation, a wallet protocol, a relayer, or a general private application
interface.  The pool address, sequence, old and new anchors, nullifier, output
commitment, asset identifier, and fee are public.

Value conservation (Theorem~\ref{thm:value-conservation}) is per spend: the
extracted output value plus the public fee equals the extracted input value,
without modular wraparound.  A sequence-wide statement about total value in
successive Merkle roots additionally needs a Merkle-binding theorem, which is
not proved here.  The initial anchor is caller supplied and this release has no
deposit or append path, so external backing and absolute solvency are also out
of scope; they would require an honestly valued initial anchor and a deposit
construction.

Theorem~\ref{thm:bcs-soundness} is a conditional classical random-oracle
argument-soundness statement in the work-normalized metric.  It makes no
quantum-random-oracle claim.  Its finite calculator checks every displayed
inequality, but it does not prove that the recorded event set completely
covers false acceptance or that the cited Johnson/MCA/circle-FRI and BCS
theorems apply to the exact mixed construction and restored-round semantics.
Those are the main unproved premises in
Assumptions~\ref{ass:pinned-finite-reduction} and
\ref{ass:boundary-restoration}.  The 124-bit Poseidon2 and 128-bit SHA entries
are policy reserves in the arithmetic ledger, not independently derived
primitive-security or second-preimage bounds; actual primitive advantages
remain symbolic assumptions.

The theorem and its 100.161-bit fractional value are scoped to the historical
q18/g37 construction, whose batch has width 29 and whose certificate uses the
conservative boundary cap \(R=32\).  Deployed V5 instead has width 19, degree
18, a batching challenge sampled from \(\mathbb K^*\), and exactly 30 S-two
public-coin rounds after the initial ensemble \(m_0\).  Its maintained theorem
proves an integer work-normalized bound at most \(2^{-100}\) only after
receiving the complete false-accept event decomposition, exact code and
literature applicability, separate-output grinding reduction, Rust transcript
correspondence, authenticated-round, PCS/Merkle, Fiat--Shamir/BCS/CMS, branch,
and actual-event correspondence premises.  No q18/g37 fractional security
number is silently transferred to V5.

Knowledge soundness and theft resistance need still stronger inputs.  The
repaired Lean theorem proves that an accepted prover execution extracting the
wrong secret is bounded by extractor failure plus a fixed-target
second-preimage event for the public nullifier.  A second Lean theorem covers a
different valid opening of the exact same fixed input leaf, reducing it to
extractor failure or a second preimage of the combined owner-and-note
commitment.  Neither theorem assumes global injectivity, which would be
impossible for a compressing hash.  A complete deployed bound still needs an
acceptance-to-extraction or simulation-extraction theorem, a defined target-
note sampling game or uniform fixed-target assumptions, computational bounds
for both commitments, a Merkle-binding reduction for an alternative leaf or
path, and the complete code-to-model connection.  The
hiding results use the stronger programmable-random-oracle experiment and do
not instantiate a standard-model pseudorandom-generator theorem, a statistical
zero-knowledge theorem, or an end-to-end deployed V5 zero-knowledge theorem.

The hiding view is exactly Definition~\ref{def:spend-complete-view}.  The
equality \(\epsilon_{\rm side}=0\) is by definition of that fixed public
\(\mathsf{Proof}\)-or-\(\mathsf{Abort}\) channel.  Local filesystem access,
the burned-nonce ledger, scheduler and process timing outside the release
event, power, thermal and memory observations, and remote prover or miner
traffic are not in the modelled view.  An implementation that exposes any of
those channels requires a new simulator and leakage bound; the present
deployment model therefore keeps rejected attempts and mining local and
private.

The hiding game assumes ideal completion before the public boundary, or an
external truncation event independent of the witness and all hidden prover
state.  The implementation instead accepts a caller-selected deadline
(3,600 seconds by default) and publishes \(\mathsf{Abort}\) when no candidate
is ready.  No wall time or selected boundary is recorded for the current
release proof.  The algebraic bound in
Lemma~\ref{lem:goodspend-product} excludes the wall-clock event and therefore is
not a bound on the total deployed Abort probability.  Applying the hiding
theorem to a concrete deadline requires either a witness-independent readiness
argument or a padded scheduling mechanism.

The privacy result also does not cover fee-payer linkage, wallet or relayer
identity, transaction origin, the account graph beyond the public statement,
or network-layer metadata.  Those are deployment privacy questions rather
than witness-hiding properties of the proof transcript.

The accepting instruction verifies and mutates state in one transaction only
after a proof account has been created, funded, uploaded, checked, and
finalized.  Those preceding transactions, their latency, fees, rent, and
storage are part of system operation but not part of the one-transaction
claim.  The mainnet run used a disposable upgradeable deployment so that
ProgramData rent could be recovered after evidence capture.  The exact code
and authority were checked through verification finality, after which a
separate close transaction removed ProgramData.  The recorded transaction is
therefore execution evidence, not a continuously available mainnet service.

The statement binds the pool public key, pool sequence and anchors, the spend
fields, and the pool's deployment domain
\(\Delta=\operatorname{SHA-256}(\mathrm{sep}\,\|\,\mathrm{pid}\,\|\,\tau)\),
where \(\mathrm{pid}\) is the runtime program id and \(\tau\) the domain tag
stored at pool initialization.  The release certificate separately binds the
frozen SBF digest.  A proof ground for one deployment domain is rejected by
any pool storing another, with a distinct error code, before proof bytes are
interpreted; cross-deployment replay against an honestly initialized pool is
therefore excluded by the statement itself rather than by operational
discipline.  One cross-cluster ambiguity remains: the holder of the program-id
keypair can redeploy the same program id on another cluster and initialize a
pool there with the same domain tag, reproducing an identical \(\Delta\).
The binding is nevertheless independently checkable: any observer can
recompute \(\Delta\) from the public program id and domain tag and confirm
which deployment a statement commits to.  Economically linked deployments
under distinct program ids or distinct domain tags cannot accept each
other's proofs.

Deployment binding is separate from nullifier derivation.  The nullifier is
\(H_{\mathsf{nul}}(k_\nu,r_{\mathrm{in}})\) under a fixed Poseidon2 domain and
does not include \(\Delta\), the note commitment, or the program id.  Reusing
the same nullifier key and input salt across deployments therefore produces
the same public nullifier and is linkable even though a proof for one stored
deployment domain is rejected by another.

The hiding theorems are conditional on the complete affine-image and
rank-coverage premise in Assumption~\ref{lem:complete-affine-image}.  The
\(\mathsf{Good}_{\mathrm{spend}}\) checker supplies release-bound evidence by reconstructing
and eliminating the relevant finite-field matrices from the frozen layout.
The independent rank checker re-derives the maps with separate field
arithmetic, while the Lean development proves the corresponding finite
algebraic hiding statements. The released q18/g37 binary is bound through
parameter CI, known-answer tests, and frozen execution evidence rather than a
universal Rust-refinement theorem.
The finalized mainnet execution establishes behaviour for one exact proof,
statement, binary, account pre-state, and transaction.  It is not a throughput
benchmark, an availability result, or an external audit, and it does not
establish resistance to denial-of-service or transaction-ordering attacks.
Proof generation for the released instance reached its configured publication
boundary; the paper does not claim a representative distribution of prover
latency or memory use.

The released transaction used a newly sampled demonstration witness rather
than a user's funded private note.  The mainnet artifact builder sampled each
secret digest from operating-system randomness, discarded candidates until the
nullifier PDA had bump 255, built the production trace, and host-verified the
proof before writing it.  The witness was not published or retained.  The
transaction therefore demonstrates private proof verification and the state
change, not custody of real value.

Finally, the selected Rust composition theorem packages component results
under successful-call, valid-input, and explicit execution/model
hypotheses.  Component C's public-run package itself carries folded-word,
coefficient, and challenge equalities.  The single hash-call equation is the
remaining equation of the Tag-67 work-verifier subtheorem, not the sole
premise of the full composition.  There is no current theorem that arbitrary
production-verifier acceptance constructs every component package or implies
the complete spend relation.
